\section*{अध्याय \ref{ch:b2:gratings} --- बहु-तरंग व्यतिकरण और जालक}

\begin{solution}{exo:b2:gratings:1}
$a = \qty{2}{\micro m}$, $\sin\theta_p = 0.275p$: \qty{16.0}{\degree}, \qty{33.4}{\degree}, \qty{55.6}{\degree};
कोई चौथी कोटि नहीं। पहली कोटि का दृश्य \qty{11.5}{\degree} से \qty{20.5}{\degree} तक:
\qty{9}{\degree}। \qty{30}{\degree} पर: $\sin\theta = 0.5 + 0.275p$: एक ओर $p = +1$ (\qty{51}{\degree}),
दूसरी ओर $p = -1$ से $-5$ ($13^\circ$, $-3^\circ$, $-19^\circ$, $-37^\circ$, $-61^\circ$)
तक।
\end{solution}

\begin{solution}{exo:b2:gratings:2}
$a = \qty{1.67}{\micro m}$। कोटि 1: $\theta = \qty{20.7}{\degree}$, $\Delta\theta = \Delta\lambda/a\cos\theta =
\qty{3.9e-4}{rad}$, \qty{0.19}{mm}; कोटि 2: $\theta = \qty{45}{\degree}$, $\qty{1.0e-3}{rad}$,
\qty{0.51}{mm}। पिक्सल \qty{50}{\micro m} या उससे कम।
\end{solution}

\begin{solution}{exo:b2:gratings:3}
$N = 12\,000$: $R = 12\,000$, $24\,000$, $36\,000$; $\Delta\lambda = 0.046$, $0.023$, \qty{0.015}{nm};
युग्मक (\qty{0.6}{nm}) और H/D जोड़ी (\qty{0.18}{nm}, जिसे $R = 3600$ चाहिए)
आसानी से अलग हो जाते हैं।
\end{solution}

\begin{solution}{exo:b2:gratings:4}
सीडी: बैंगनी \qty{14.5}{\degree}, लाल \qty{25.9}{\degree}; लाल कोटि 2 तक है, बैंगनी
कोटि 4 तक: अर्थात् दो (कुछ चढ़ते हुए) इंद्रधनुष। DVD: बैंगनी \qty{32.7}{\degree}, लाल
\qty{71}{\degree}; लाल की केवल पहली कोटि है: एक इंद्रधनुष, चौड़ा।
\end{solution}

\begin{solution}{exo:b2:gratings:5}
(क) शून्य $\varphi = 2\pi m/N$ पर, $m = 1, \dots, N - 1$: अर्थात् $N - 1$ शून्य, और उनके बीच $N - 2$
अधिकतम। (ख) $N\varphi/2 = 3\pi/2$ पर, $I \approx I_0/\sin^2(3\pi/2N) \approx I_0(2N/3\pi)^2$:
अर्थात् $4/9\pi^2 = 0.045$ गुना $N^2I_0$। (ग) $N = 3$: $\varphi = \pi$ पर ऊँचाई $I_0$ का कोई
अकेला \omterm{thm:b2:gratings:nwaves}{गौण अधिकतम}, यानी मुख्य $9I_0$ का नवाँ भाग; $N = 6$: कुछ प्रतिशत के चार
\omterm{thm:b2:gratings:nwaves}{गौण अधिकतम}। (घ) $\sin^2x/x^2$ का केंद्रीय लोब
ऊर्जा का $90\%$ रखता है।
\end{solution}

\begin{solution}{exo:b2:gratings:6}
(क) $p > \lambda_{\min}/(\lambda_{\max} - \lambda_{\min}) = 1.33$: अर्थात् $p = 2$ से। (ख) $\lambda/p$। (ग)
\qty{25}{nm}; और जालक के साथ आड़ा रखा कोई प्रिज़्म कोटियों को एक के ऊपर एक
चुन देता है। (घ) $\Delta\lambda = \lambda^2/2L$, अर्थात् किसी नैनोमीटर का अंश --- किसी जालक वाले से
सैकड़ों गुना छोटा।
\end{solution}

\begin{solution}{exo:b2:gratings:7}
(क) $\sin\beta = \lambda/2a = 0.3$, $\beta = \qty{17.5}{\degree}$। (ख) अकेला फलक अपनी दर्पणवत् दिशा
पर केंद्रित किसी चौड़े लोब में विवर्तित करता है; और जो कोटियाँ
उसके भीतर पड़ती हैं उन्हें प्रकाश मिलता है। (ग) कोटि 2 (उसी कोण पर $2 \times \qty{500}{nm} = \qty{1}{\micro m}$)।
(घ) $R = 2W\sin\beta/\lambda = 3.6 \times 10^5$।
\end{solution}

\begin{solution}{exo:b2:gratings:8}
(क) $\dd\lambda/\dd x = a\cos\theta/pf = \qty{3.1}{nm/mm}$: \qty{50}{\micro m} \qty{0.16}{nm} है। (ख)
$R = 30\,000$: \qty{0.018}{nm}। (ग) झिर्री, नौ गुना से; और बराबर $s = \qty{6}{\micro m}$ पर।
(घ) झिर्री पर विवर्तन और कोई भूखा संसूचक।
\end{solution}

\begin{solution}{exo:b2:gratings:9}
(क) $p = 2L/\lambda = 16\,700$; \omterm{prop:b2:gratings:fabryperot}{मुक्त वर्णक्रमी परास} $c/2L = \qty{30}{GHz}$, $\lambda^2/2L = \qty{0.036}{nm}$। और (ख) $\mathcal F
= \pi\sqrt{0.95}/0.05 = 61$; \qty{0.5}{GHz}, \qty{0.6}{pm}। (ग) $p\mathcal F = 10^6$, अर्थात्
\qty{10}{cm} के जालक का चार गुना। (घ) उसका \omterm{prop:b2:gratings:fabryperot}{मुक्त वर्णक्रमी परास} किसी नैनोमीटर का
अंश है: अतः कोटियाँ चढ़ जाती हैं जब तक प्रकाश उनमें से किसी एक तक
पूर्व-छाना न जाए।
\end{solution}

\begin{solution}{exo:b2:gratings:10}
(क) $|\sin\theta - \sin\theta_0| \le 2$। (ख) चरम किरणों के पथ में
$W(\sin\theta - \sin\theta_0)$ का अंतर है, अधिक-से-अधिक $2W$। (ग) $2 \times 0.3/5 \times 10^{-7} = 1.2 \times 10^6$; $c/R
= \qty{250}{m/s}$। (घ) किसी रेखा का केंद्र उसकी चौड़ाई के सौवें भाग तक
जान लिया जाता है, और हज़ार रेखाएँ $\sqrt{1000}$ से माध्य घटा देती हैं: अर्थात् मीटर प्रति
सेकंड, वीरतापूर्ण स्थिरता के साथ।
\end{solution}

\begin{solution}{exo:b2:gratings:11}
(क) निचले तल से परावर्तित तरंग $2d\sin\theta$ अधिक चलती है।
(ख) $\sin\theta = 0.273p$: \qty{15.8}{\degree}, \qty{33.1}{\degree}, \qty{55.0}{\degree}। (ग) $\lambda \gg 2d$:
कोई कोण समीकरण नहीं मानता; और काँच के कोई नियमित तल नहीं हैं, अतः
केवल विसरित वलय --- वह अक्रिस्टली है। (घ) $\sin\theta \le 1$।
\end{solution}

\begin{solution}{exo:b2:gratings:12}
(क) पड़ोसियों की कला पथ से $2\pi a\sin\theta/\lambda$ और निवेश से $-\psi$
भिन्न है; और \omterm{thm:b2:gratings:nwaves}{मुख्य अधिकतम} $\varphi = 0$ पर। (ख) पहला शून्य
$\Delta\varphi = 2\pi/N$ पर: $\Delta\theta = \lambda/Na\cos\theta$। (ग) $\varphi = \pm2\pi$ पर कोई दूसरा \omterm{thm:b2:gratings:nwaves}{मुख्य अधिकतम}
$|\sin\theta \pm \lambda/a| \le 1$ माँगता है: और यह हर मोड़ने वाले कोण के लिए तभी असंभव है
जब $a \le \lambda/2$ हो। (घ) $Na = \qty{0.96}{m}$: \qty{1.8}{\degree}; $50$ क़दम,
\qty{50}{\micro s}।
\end{solution}

\begin{solution}{pb:b2:gratings:1}
\textbf{1.} $a = \qty{0.833}{\micro m}$; $N = 144\,000$।

\textbf{2.} $\sin\theta = \sin20^\circ - \lambda/a$ (कोटि अभिलंब की दूसरी ओर पड़ती
है): \qty{18.5}{\degree} पर \qty{550}{nm}, \qty{7.9}{\degree} पर \qty{400}{nm}, \qty{29.9}{\degree} पर
\qty{700}{nm}: अर्थात् वर्णक्रम के \qty{22}{\degree}।

\textbf{3.} $1/a\cos\theta = \qty{1.27e-3}{rad/nm}$; $f_2\times$: \qty{1.27}{mm/nm}, अर्थात्
\qty{0.79}{nm/mm}, और प्रति पिक्सल \qty{0.012}{nm}।

\textbf{4.} $R = N = 144\,000$: \qty{0.0038}{nm}, अर्थात् किसी पिक्सल का तिहाई।

\textbf{5.} दूसरी कोटि का \qty{400}{nm} वहाँ उतरता है जहाँ पहली कोटि का \qty{800}{nm}
उतरता, अर्थात् लाल के परे: दृश्य में कोई चढ़ाव नहीं; चढ़ाव
\qty{350}{nm} पर शुरू होता, जिसे कोई काँच का छन्ना हटा देता है।

\textbf{6.} \qty{22}{\degree} $\times$ \qty{1}{m} $= \qty{38}{cm}$: कोई संसूचक एक फाँक ढकता है;
और उसे चुनने के लिए जालक घुमाया जाता है।

\textbf{7.} $\lambda = 2a\sin17^\circ = \qty{490}{nm}$।

\textbf{8.} $0.8 \times 0.4 = 32\%$।

\textbf{9.} कोण \qty{26.5}{\degree}, \qty{14.0}{\degree}, \qty{10.3}{\degree}: \qty{550}{nm} की जगह से $+14$, $-8$ और
\qty{-14}{cm}।

\textbf{10.} \qty{0.18}{nm} $15$ पिक्सल है: अलग हुई।

\textbf{11.} $50$ पिक्सल की दूरी पर, हर एक $4$ पिक्सल चौड़ी।

\textbf{12.} \qty{100}{\micro m}, $6.7$ पिक्सल, \qty{0.08}{nm} --- अर्थात् जालक की सीमा का
बीस गुना: झिर्री-सीमित।

\textbf{13.} प्रतिबिंब \qty{97}{\micro m} चौड़ा है: कोई \qty{50}{\micro m} की झिर्री
आधा प्रकाश जाने देती है।

\textbf{14.} उतनी ही तीखाई के लिए कम प्रकाश: फ़ोटॉन रव बढ़ जाता है।

\textbf{15.} $\Delta\lambda = \lambda v/c = \qty{0.066}{nm}$, $5.5$ पिक्सल।

\textbf{16.} पृथ्वी का वेग वर्ष भर हर रेखा को $\pm5.5$ पिक्सल तक
खिसका देता है; और वह पंचांगों से \qty{100}{m/s} के लिए ज़रूरी $0.3\%$ से कहीं
बेहतर ज्ञात है।

\textbf{17.} \qty{2.2e-5}{nm}, $0.002$ पिक्सल --- रेखा की चौड़ाई से कहीं नीचे। किसी
\qty{7}{px} रेखा के सौवें भाग तक कोई केंद्रक \qty{0.07}{px} है; और हज़ार
रेखाएँ $\sqrt{1000}$ देती हैं: $0.002$ पिक्सल --- अर्थात् संकेत, बस मुश्किल से। असली
यंत्र और विभेदन-शक्ति तथा स्थिरता जोड़ते हैं।

\textbf{18.} $\Delta\sin\theta = (\lambda/a) \times 10^{-5} = 6.6 \times 10^{-6}$, $\Delta\theta = \qty{7e-6}{rad}$,
\qty{7}{\micro m}: अर्थात् प्रति केल्विन आधा पिक्सल --- और \qty{10}{m/s} के लिए मिलीकेल्विन।

\textbf{19.} उसकी रेखाएँ, उसी क्षण उन्हीं पिक्सलों पर दर्ज,
तरंगदैर्घ्य का पैमाना देती हैं और हर बहाव का पीछा करती हैं।

\textbf{20.} $a = \qty{3.33}{\micro m}$, $N = 36\,000$, वही कोण (\qty{18.5}{\degree}), $R = 4N
= 144\,000$, वही परिक्षेपण, \omterm{prop:b2:gratings:fabryperot}{मुक्त वर्णक्रमी परास} $\lambda/4 = \qty{137}{nm}$: अर्थात् विभेदन-शक्ति और
परिक्षेपण केवल $W$ और कोण पर निर्भर हैं; और \omterm{prop:b2:gratings:fabryperot}{मुक्त वर्णक्रमी परास}
सिकुड़ गया।

\textbf{21.} \omterm{prop:b2:gratings:fabryperot}{मुक्त वर्णक्रमी परास} $\lambda^2/2L = \qty{0.15}{nm}$ ($12$ पिक्सल); $\mathcal F = 30$; चोटियाँ
\qty{0.005}{nm} चौड़ी: अर्थात् संसूचक भर तीखी रेखाओं की कोई कंघी, कोई
तरंगदैर्घ्य-पैमाना।

\textbf{22.} \qty{1.3e-3}{rad/nm} के सामने $2 \times 10^{-4}$: छह गुना कम, और किसी
प्रिज़्म की विभेदन-शक्ति ($b\,\dd n/\dd\lambda \sim 10^4$) तथा अरैखिक पैमाना भी
हारते हैं; और जालक पराबैंगनी में भी काम करते हैं।

\textbf{23.} $R \le 2W/\lambda$: \qty{500}{nm} पर मीटर भर के लिए $4 \times 10^6$ --- चौड़ाई और
कोण ही ऊपर जाने के रास्ते हैं।

\textbf{24.} $10^4 \times 0.012 \times 0.32 = 38$ फ़ोटॉन प्रति पिक्सल प्रति सेकंड; \qty{100}{s} में
$3800$: संकेत-रव अनुपात $\approx 60$।

\textbf{25.} परिक्षेपण, जिसे $p/a\cos\theta$ और कैमरे की फ़ोकस दूरी तय करती है;
विभेदन-शक्ति, जिसे जगमगाई चौड़ाई और कोण ($pN$) तय करते हैं;
और \omterm{prop:b2:gratings:fabryperot}{मुक्त वर्णक्रमी परास}, जिसे कोटि ($\lambda/p$) तय करती है।
\end{solution}
